\section{模型检验与效率口径对照}
\subsection{主结果的物理和费用核验}
主模型已有验证记录包含储电量递推、容量与功率上下界、充放电互斥、跨日连续性、紧急缺口、逐段结算、日汇总及结果表反读。第一问共144段，每个年度结果集分别包含334天、48096段，不能将48096写成全部四问合计。

第一问记录的能量平衡最大残差约$2.27\times10^{-13}$ kWh，储电量递推最大残差约$7.28\times10^{-12}$ kWh。第三问与4-3均无执行充放电重叠，储电量递推最大误差约$4.55\times10^{-12}$ kWh，实际费用分别复算为14022396.98元和15061735.43元。这里引用既有记录，不将此次论文整理描述为重新运行全部求解器。

费用复算能确认给定调度的账目正确，却不能单独证明优化目标正确。第\ref{sec:q43objective}节所示4-3目标附加项不会因实际费用核验通过而消失；该问题需将原目标与修正目标进行独立求解对照后才能量化影响。

\subsection{信息边界与结算解释}
已有程序通过扰动未来实际值及尚未发布预报检查当前预测与决策是否改变；这类检查针对执行阶段的前视读取。它不能证明原预测参数在历史研发中从未参考测试年，因此本文将结论限定为给定参数下的历史滚动回测，不称为完全独立的外部验证。

取消购电的退款规则会影响实际费用。按原主调度不变、取消不退原费且仍另付50\%违约费的解释，第三问费用为14172567.11元，4-3为15234936.02元。这是同一调度的结算敏感性，不是另一合同下重新优化的最优费用。正文主结果始终采用式\eqref{eq:q3settle}，不得将两种结算混用。

\subsection{储能效率口径对结果的影响}\label{sec:efficiency}
题面仅给出“充放电效率90\%”，没有进一步区分单向与往返。主模型按单向各0.9建立；已有补充材料\cite{efficiency}另采用往返0.9并对称拆分损耗。两者的物理递推分别为
\begin{align}
\text{主口径：}\quad&E_t=E_{t-1}+0.9c_t-b_t/0.9,\\
\text{对照口径：}\quad&E_t=E_{t-1}+\eta c_t-b_t/\eta,\quad\eta=\sqrt{0.9}.
\end{align}
对称拆分是对照模型的明确选择，不意味着相同往返效率下所有单向拆分都会产生完全相同的受限调度。

\begin{table}[htbp]\centering\small
\caption{已有两种效率口径总费用对照（元）}\label{tab:efficiency}
\begin{tabular}{lrrrr}\toprule
结果集 & 往返0.81 & 往返0.90 & 差额 & 变化率\\\midrule
问题一 & 35126.95 & 33801.50 & -1325.45 & -3.77\%\\
问题二 & 14389135.08 & 13910066.83 & -479068.25 & -3.33\%\\
问题三 & 14022396.98 & 13558297.17 & -464099.81 & -3.31\%\\
问题四-2 & 15163043.89 & 14684028.55 & -479015.34 & -3.16\%\\
问题四-3 & 15061735.43 & 14579714.65 & -482020.78 & -3.20\%\\
\bottomrule\end{tabular}\end{table}

表\ref{tab:efficiency}逐项列示补充材料已经报告的总费用。对照口径下五组费用下降约3.16\%—3.77\%，说明效率解释对费用绝对值有可观影响，应在报告结果时声明，不能视为可忽略的数值误差。降低损耗会改变可用输出电量与库存安排，因此这种变化在机理上合理。

该表只有五组主结果的总费用，尚不足以证明所有更新频率排序、均价与联合场景排序在两种效率下都保持不变。本文也不以总费用差额断言某一独立套利或光伏消纳效应。效率对照采用补充材料的既有记录，本次未读取各问效率分支或重新核验其逐段输出，故不新增“两口径全部逐段核验通过”的结论。若用于最终定稿，应将相关论断限制在现有表格可支持的范围。

\subsection{图表与解释一致性}
所有主结果图取自新图包\cite{figures}，保持统一字体、配色与单位，并保留来源对应关系。差值图明确相减方向，费用图区分元和万元，紧急时段标记不冒充电量；热力图不截断极端值。第四问跨情形费用比较同时伴随模型条件变化，属于描述性比较。图区分了“变化是什么”与“变化由何导致”，避免用视觉相似性或相关性替代因果证据。

\section{模型评价}
\subsection{优点与适用范围}
共同的储能约束使五组结果能够以一致的物理单位解释；历史场景和滚动执行将预测风险、购电费用与库存价值联系起来。第三问通过节点共享变量表达信息限制，能够在不预知未来实际值的条件下利用预报更新。费用分项与同价结算对照使节省来源和比较口径可追溯，指定日期图表又提供了对全年汇总的具体解释。

\subsection{局限与改进方向}
首先，4-3当前目标含合同费用之外的附加项，现有结果不能证明针对合同总费用的最优性。优先改进应是核对目标系数与合同解释，再决定是否重算，不应直接将该项事后包装成风险偏好。

其次，经验场景来自有限历史窗口，小型二叉树和次日确定性近似无法覆盖全部随机结构。窗口、预测器及融合收缩参数虽在执行时固定，其跨年适用性仍需独立数据验证。第二问启动期改善和第三问18点微小差异，均不宜被放大为普遍规律。

最后，效率、功率作用位置、取消退款和定价时刻均存在明确建模解释；储能寿命、设备退化和售电机制不在现有主模型内。后续若增加这些因素，需要同步修改状态或成本，而非沿用现有费用直接外推。
